
\begin{abstract}
   With the rapid growth of urban instant delivery, scheduling for new energy logistics fleets faces increasing challenges from dynamic task arrivals, limited vehicle resources, and complex charging constraints. This paper presents a dynamic collaborative scheduling system based on graph modeling and intelligent decision methods. The urban road network is represented as a weighted graph, and a discrete-time simulation framework is built by considering battery limits, load capacity, task deadlines, and charging station congestion. 
   For path planning and reachability analysis, the system applies Dijkstra, A*, and RRT, while at the scheduling level it implements several heuristic strategies, including nearest-task-first, maximum-weight-first, earliest-deadline-first and Q-learning-based hyper-heuristic. 
   In addition, an offline MILP module under full-information settings is developed for small-scale cases to generate near-optimal solutions and provide a baseline for comparison with online scheduling strategies. 
   The system follows a “simulation engine + backend + visualization frontend” architecture. Experimental show that it can effectively capture the key constraints and operating process of new energy logistics delivery, and that different strategies exhibit clear differences in task completion rate, total reward, path length, and charging load. 
   Overall, this work develops a research-relevant platform for dynamic scheduling of new energy logistics fleets, providing a practical basis for future studies on multi-agent coordination, exact optimization, and more complex urban delivery scenarios.
\end{abstract}